%\documentclass[12pt]{article}
%\usepackage{amsmath}
%\usepackage{amssymb}
%\usepackage{geometry}
%\geometry{a4paper, margin=1in}

%\title{GATE BM 2025 --- Theory / Formulae Used}
%\author{}
%\date{}

%\begin{document}
%\maketitle

\noindent This document lists only the \textbf{general theory and formulae}
used to solve each question (no numerical substitution).
\begin{enumerate}[label=\thesubsection.\arabic*,ref=\thesubsection.\theenumi]
% =========================================================
%\section*{1. Geometric Progression of Frequencies}
\item Geometric Progression of Frequencies

If each successive term is $r$ times the previous term, then after
$k$ steps from a reference term $f_0$,

\begin{align}
\frac{f_k}{f_0}=r^{k}. \label{GP}
\end{align}


%\end{enumerate}
% =========================================================
\item Self-Similar (Fractal) Curve Length

If, at every iteration, each segment is replaced by $m$ segments each
of length $\dfrac{1}{p}$ of the original segment, the length scales as

\begin{align}
L_1=L_0\left(\frac{m}{p}\right),
\qquad
L_n=L_0\left(\frac{m}{p}\right)^{n}. \label{iteration}
\end{align}

% =========================================================
\item Triangle Inequality

A triangle can be formed from three lengths $l_1,l_2,l_3$ if and only if
the sum of any two sides exceeds the third:

\begin{align}
l_1+l_2>l_3,\qquad
l_2+l_3>l_1,\qquad
l_3+l_1>l_2.
\end{align}

% =========================================================
\item Notch (Band-Stop) Filter Response

A notch filter strongly attenuates a narrow band of frequencies around
a rejected frequency $\omega_0$, while passing frequencies away from
it:

\begin{align}
|H(j\omega_0)|\approx 0,
\qquad
|H(j\omega)|\ \text{high for}\ \omega<\omega_0\ \text{and}\ \omega>\omega_0. \label{freq}
\end{align}

% =========================================================
\item Rolle's Theorem \label{rolle's}

If $f$ is continuous on $[a,b]$, differentiable on $(a,b)$, and
$f(a)=f(b)$, then

\begin{align}
\exists\, c\in(a,b)\ \text{such that}\ f'(c)=0.
\end{align}

Contrapositive form: if $f'(x)\neq 0$ for all $x$, then
$f(a)\neq f(b)$ for $a\neq b$.

% =========================================================
\item Standard Second-Order System / Closed-Loop Transfer Function

For unity negative feedback with forward transfer function $G(s)$,

\begin{align}
\frac{C(s)}{R(s)}=\frac{G(s)}{1+G(s)}. \label{neg}
\end{align}

Comparing the characteristic equation with the standard second-order
form,

\begin{align}
s^{2}+2\zeta\omega_n s+\omega_n^{2}=0,
\end{align}

gives the natural frequency $\omega_n$ and damping ratio $\zeta$.
The system is:

\begin{itemize} \label{compare}
\item overdamped if $\zeta>1$,
\item critically damped if $\zeta=1$,
\item underdamped if $0<\zeta<1$,
\item undamped (sustained oscillations) if $\zeta=0$.
\end{itemize}

For $0<\zeta<1$, the impulse response has the form

\begin{align}
c(t)=A\,e^{-\zeta\omega_n t}\sin\!\left(\omega_n\sqrt{1-\zeta^{2}}\;t\right),
\qquad t\ge 0. \label{inverse}
\end{align}

% =========================================================
\item Eigenvalues via the Characteristic Equation

For a matrix $A$, the eigenvalues $\lambda$ satisfy

\begin{align}
\det(A-\lambda I)=0.
\end{align}

For a $2\times2$ rotation matrix

\begin{align}
A=
\begin{pmatrix}
\cos\theta & -\sin\theta\\
\sin\theta & \cos\theta
\end{pmatrix},
\end{align}

using $\sin^2\theta+\cos^2\theta=1$ and Euler's formula
$e^{\pm i\theta}=\cos\theta\pm i\sin\theta$, the eigenvalues reduce to

\begin{align}
\lambda=\cos\theta\pm i\sin\theta=e^{\pm i\theta}.
\end{align}

% =========================================================
\item Definition of the Derivative (First Principles)

\begin{align}
f'(x_0)=\lim_{\Delta x\to 0}\frac{f(x_0+\Delta x)-f(x_0)}{\Delta x}. \label{derivative}
\end{align}

% =========================================================
\item Mean and Median of a Dataset

For $n$ ordered values $x_1\le x_2\le\dots\le x_n$:

\begin{align}
\text{Mean}=\frac{1}{n}\sum_{i=1}^{n}x_i. \label{mean}
\end{align}

If $n$ is odd, the median is the middle term,

\begin{align}
\text{Median}=x_{\frac{n+1}{2}}. \label{median}
\end{align}

% =========================================================
\item Nyquist Sampling Theorem \label {sampling}

A band-limited signal with maximum frequency component $f_{\max}$
must be sampled at a rate

\begin{align}
f_s\ge 2f_{\max}
\end{align}

to allow accurate (aliasing-free) reconstruction. For an angular
frequency component $\omega$, the corresponding cyclic frequency is

\begin{align}
f=\frac{\omega}{2\pi}.
\end{align}

% =========================================================
\item Average Power in an AC Circuit

\begin{align}
P_{\text{avg}}=V_{\text{rms}}I_{\text{rms}}\cos\theta, \label{power}
\end{align}

where

\begin{align}
V_{\text{rms}}=\frac{V_m}{\sqrt2},
\qquad
I_{\text{rms}}=\frac{I_m}{\sqrt2},
\qquad
\theta=\phi_v-\phi_i.
\end{align}

% =========================================================
\item Mean and Variance of a Binomial Distribution

For $X\sim B(N,p)$,

\begin{align}
E(X)=Np,
\qquad
\operatorname{Var}(X)=Np(1-p). \label{variance}
\end{align}

% =========================================================
\item Solving a Linear ODE via the Laplace Transform

For $y''(x)-k^2y(x)=0$ with initial conditions $y(0)$, $y'(0)$, take
the Laplace transform using

\begin{align}
\mathcal{L}\{y''(x)\}=s^{2}Y(s)-s\,y(0)-y'(0). \label{laplace}
\end{align}

This gives an algebraic equation for $Y(s)$, which is solved by
partial fractions,

\begin{align}
Y(s)=\frac{A}{s-a}+\frac{B}{s+a}, \label{partial}
\end{align}

and inverted using

\begin{align}
\mathcal{L}^{-1}\left\{\frac{1}{s-a}\right\}=e^{at}.
\end{align}

% =========================================================
\item Impedance of R--C Networks (Phasor Domain) \label{impedence deri}



\text{Symbolic Impedance Derivation}

The impedance of the capacitor in terms of angular frequency $\omega$ is:
\begin{align}
Z_C(\omega) = \frac{1}{j\omega C} = -\frac{j}{\omega C}
\end{align}

The parallel combination of resistor $R_2$ and capacitor $Z_C(\omega)$ gives $Z_p(\omega)$:
\begin{align}
Z_p(\omega) &= \frac{R_2 \cdot Z_C(\omega)}{R_2 + Z_C(\omega)} = \frac{R_2 \left(-\frac{j}{\omega C}\right)}{R_2 - \frac{j}{\omega C}}\\
& = \frac{-j R_2}{\omega R_2 C - j}
\end{align}

Multiplying numerator and denominator by the complex conjugate $(\omega R_2 C + j)$:
\begin{align}
Z_p(\omega) &= \frac{-j R_2 (\omega R_2 C + j)}{(\omega R_2 C - j)(\omega R_2 C + j)}\\
&= \frac{R_2 - j \omega R_2^2 C}{1 + (\omega R_2 C)^2}
\end{align}

The total input impedance $Z_{\text{in}}(\omega)$ in series with $R_1$ is:
\begin{align}
Z_{\text{in}}(\omega) &= R_1 + Z_p(\omega)\\
&= \left( R_1 + \frac{R_2}{1 + (\omega R_2 C)^2} \right) - j \left( \frac{\omega R_2^2 C}{1 + (\omega R_2 C)^2} \right)
\end{align}

\vspace{0.5em}
\text{Magnitude and Phase Expressions in terms of $\omega$}

\text{1. Magnitude Spectrum $|Z_{\text{in}}(\omega)|$:}
\begin{align}
|Z_{\text{in}}(\omega)| = \sqrt{ \left[ R_1 + \frac{R_2}{1 + (\omega R_2 C)^2} \right]^2 + \left[ \frac{\omega R_2^2 C}{1 + (\omega R_2 C)^2} \right]^2 }
\end{align}

Simplifying the algebraic expression:
\begin{align}
|Z_{\text{in}}(\omega)| = \sqrt{ R_1^2 + \frac{2 R_1 R_2 + R_2^2 + (\omega R_1 R_2 C)^2}{1 + (\omega R_2 C)^2} }
\end{align}

\text{2. Phase Spectrum $\angle Z_{\text{in}}(\omega)$:}
\begin{align}
\angle Z_{\text{in}}(\omega) &= -\tan^{-1}\left( \frac{\frac{\omega R_2^2 C}{1 + (\omega R_2 C)^2}}{R_1 + \frac{R_2}{1 + (\omega R_2 C)^2}} \right)\\
& = -\tan^{-1}\left( \frac{\omega R_2^2 C}{R_1(1 + (\omega R_2 C)^2) + R_2} \right)
\end{align}

% =========================================================
\item Complex Logarithm \label{log}

For real $x$, using Euler's identity for $-1$,


\begin{align}
y = \log_e(-e^x) = \ln(-e^x)
\end{align}

Using logarithmic properties, we rewrite the term inside the logarithm:

\begin{align}
-e^x = (-1) \cdot e^x
\end{align}

Applying Euler's identity, $-1$ can be expressed in complex exponential form as:

\begin{align}
-1 = e^{i(2k+1)\pi}, \quad \text{for } k \in \mathbb{Z}
\end{align}

By solving:

\begin{align}
-e^x = e^{x} \cdot e^{i(2k+1)\pi} = e^{x + i(2k+1)\pi} \label{e^x}
\end{align}




% =========================================================
\item Linearity and Causality of a System

\textbf{Linearity} (superposition): a system is linear if, for inputs
$x_1,x_2$ and scalars $a,b$,

\begin{align}
T\{a\,x_1(t)+b\,x_2(t)\}=a\,T\{x_1(t)\}+b\,T\{x_2(t)\}.
\end{align}

\textbf{Causality}: a system is causal if the output at time $t_0$
depends only on input values for $t\le t_0$; if it depends on future
input values ($t>t_0$), it is noncausal.

% =========================================================
\item Linear (Discrete) Convolution

For an LTI system with impulse response $h[n]$ of length $M$ and
input $x[n]$ of length $N$,

\begin{align}
y[n]=h[n]*x[n]=\sum_{k}h[k]\,x[n-k],
\end{align}

with output length

\begin{align}
L=M+N-1. \label{convolution}
\end{align}

This can be written as a matrix--vector product using a lower
triangular Toeplitz matrix \myvec{H} built from shifted copies of
$h[n]$:

\begin{align}
\myvec{y}=\myvec{H},\myvec{x}.
\end{align}

% =========================================================
\item $n$-bit DAC Output

For an $n$-bit DAC with full-scale output voltage $V_{\text{FS}}$ and
digital input equivalent decimal value $D$,

\begin{align}
V_{\text{out}}=V_{\text{FS}}\left(\frac{D}{2^{n}-1}\right). \label{output}
\end{align}

% =========================================================
\item Inverting Summing Amplifier

For an op-amp inverting summing amplifier with feedback resistor
$R_f$, input resistors $R_1,R_2,R_3,\dots$ and corresponding input
voltages $V_1,V_2,V_3,\dots$, the output voltage is

\begin{align}
V_{\text{out}}=-\left(\frac{R_f}{R_1}V_1+\frac{R_f}{R_2}V_2+\frac{R_f}{R_3}V_3+\cdots\right), \label{amplifier}
\end{align}

valid provided $V_{\text{out}}$ stays within the saturation limits,

\begin{align}
-V_{\text{sat}}\le V_{\text{out}}\le +V_{\text{sat}}.
\end{align}

\item Derivation

Let the inverting input node (where $R_1,R_2,R_3,R_f$ meet) be node
$A$, and let the non-inverting input be grounded.

\textbf{Ideal op-amp assumptions:}

\begin{align}
\text{(i)}\quad &\text{No current flows into either input terminal.}\\
\text{(ii)}\quad &\text{The open-loop gain is infinite, so the two
inputs are forced to the same potential (virtual short).}
\end{align}

Since the non-inverting input is at $0\,\text{V}$ (grounded), the
virtual short forces node $A$ to also be at $0\,\text{V}$:

\begin{align}
V_A=0
\qquad(\text{virtual ground}).
\end{align}

Applying Kirchhoff's Current Law (KCL) at node $A$, the sum of
currents flowing into the node through $R_1,R_2,R_3$ must equal the
current flowing out through the feedback resistor $R_f$ (since no
current enters the op-amp input):

\begin{align}
\frac{V_1-V_A}{R_1}+\frac{V_2-V_A}{R_2}+\frac{V_3-V_A}{R_3}
=\frac{V_A-V_{\text{out}}}{R_f}.
\end{align}

Substituting $V_A=0$:

\begin{align}
\frac{V_1}{R_1}+\frac{V_2}{R_2}+\frac{V_3}{R_3}
=\frac{-V_{\text{out}}}{R_f}.
\end{align}

Solving for $V_{\text{out}}$:

\begin{align}
V_{\text{out}}
=-R_f\left(\frac{V_1}{R_1}+\frac{V_2}{R_2}+\frac{V_3}{R_3}\right)
=-\left(\frac{R_f}{R_1}V_1+\frac{R_f}{R_2}V_2+\frac{R_f}{R_3}V_3\right),
\end{align}

which recovers equation (19.1). The same derivation extends to any
number of input branches $R_i,V_i$:

\begin{align}
V_{\text{out}}=-\sum_{i}\frac{R_f}{R_i}V_i.
\end{align}

% =========================================================
\item Wien Bridge Oscillator Frequency

For a Wien bridge oscillator with equal frequency-selective resistors
$R$ and capacitors $C$, the oscillation frequency is

\begin{align}
f_0=\frac{1}{2\pi RC}. \label{wien}
\end{align}

% =========================================================
\item Series RLC Circuit --- Resonance \label{resonance}

Kirchhoff's Voltage Law for a series $RLC$ circuit driven by
$v(t)=V_m\cos(\omega t)$:

\begin{align}
R\,i(t)+L\frac{di(t)}{dt}+\frac{1}{C}\int i(t)\,dt=v(t).
\end{align}

In the phasor domain, the impedance is

\begin{align}
Z(\omega)=R+j\left(\omega L-\frac{1}{\omega C}\right)\\
\qquad
|Z(\omega)|=\sqrt{R^{2}+\left(\omega L-\frac{1}{\omega C}\right)^{2}}.
\end{align}

Converting to complex phasor notation $V(t) = V_{\text{rms}} \sqrt{2} e^{j\omega t}$ and $I(t) = I_{\text{rms}} \sqrt{2} e^{j(\omega t - \phi)}$:
\begin{align}
Z(\omega) = \frac{V(\omega)}{I(\omega)} = R + j\omega L + \frac{1}{j\omega C} = R + j\left(\omega L - \frac{1}{\omega C}\right)
\end{align}

Resonance occurs when the reactive part vanishes:

\begin{align}
\omega_0 L=\frac{1}{\omega_0 C}
\;\Longrightarrow\;
\omega_0=\frac{1}{\sqrt{LC}},
\qquad
f_0=\frac{\omega_0}{2\pi}.
\end{align}

At resonance, $Z(\omega_0)=R$, so

\begin{align}
I_{\text{rms}}(\omega_0)=\frac{V_{\text{rms}}}{R}. \label{res}
\end{align}

% =========================================================
\item Newton--Raphson Method

For a root of $f(x)=0$, the iterative update starting from $x_n$ is

\begin{align}
x_{n+1}=x_n-\frac{f(x_n)}{f'(x_n)}. \label{raphson}
\end{align}

% =========================================================
\item Logarithmic Decay Model \label {decay}

For a quantity that decays logarithmically with time,

\begin{align}
\sigma(t)=\sigma_0-k\ln(t),\qquad t\ge 1,
\end{align}

the difference between two times $t_1,t_2$ eliminates the unknown
constant $\sigma_0$:

\begin{align}
\sigma(t_2)-\sigma(t_1)=-k\left[\ln(t_2)-\ln(t_1)\right]
=-k\ln\!\left(\frac{t_2}{t_1}\right).
\end{align}

% =========================================================
\item Circular Elimination (Josephus-type) Counting \label{count}

Given $N$ people arranged in a circle and a fixed direction of
counting, the $k^{\text{th}}$ person from a reference position $P$ in
that direction is located by counting $k$ steps around the circle;
eliminated members are removed from the circle before subsequent
counts are made from the new reference position.

% =========================================================
\item Percentage / Proportion of a Total

For a subset with combined value $S$ out of a grand total $T$,

\begin{align}
\text{Percentage}=\frac{S}{T}\times 100\%.
\end{align}

% =========================================================
\item Piecewise (Slab-Based) Fee Structure with Tax

For a fee $F$ exceeding a threshold $F_0$, with base overhead $O_0$
and marginal overhead rate $r_1$ beyond the threshold, plus a Goods
and Services Tax rate $r_2$ on the (fee + overhead):

\begin{align}
O=O_0+r_1\left(F-F_0\right),
\end{align}

\begin{align}
T_{\text{pre-tax}}=F+O,
\end{align}

\begin{align}
T_{\text{total}}=(1+r_2)\,T_{\text{pre-tax}}.
\end{align}

Given a maximum payable $T_{\text{total}}$, this is solved for the
maximum permissible $F$.

% =========================================================
\item Gaussian (Normal) Function

A random variable $X$ is said to be Gaussian (normally distributed)
with mean $\mu$ and variance $\sigma^{2}$, written
$X\sim\mathcal{N}(\mu,\sigma^{2})$, if its probability density
function is

\begin{align} \label{gaussian}
f(x)=\frac{1}{\sigma\sqrt{2\pi}}\,
\exp\!\left(-\frac{(x-\mu)^{2}}{2\sigma^{2}}\right),
\qquad -\infty<x<\infty.
\end{align}

The \text{standard normal} form, with $\mu=0,\ \sigma=1$, is

\begin{align}
\phi(z)=\frac{1}{\sqrt{2\pi}}\,e^{-z^{2}/2},
\qquad z=\frac{x-\mu}{\sigma}.
\end{align}

The function is symmetric about $x=\mu$, and satisfies

\begin{align}
\int_{-\infty}^{\infty}f(x)\,dx=1,
\qquad
E[X]=\mu,
\qquad
\operatorname{Var}(X)=\sigma^{2}.
\end{align}

% =========================================================
\item Central Limit Theorem (CLT)

Let $X_1,X_2,\dots,X_n$ be independent and identically distributed
(i.i.d.) random variables, each with finite mean $\mu$ and finite
variance $\sigma^{2}$ (the underlying distribution need not itself be
Gaussian). Define the sample mean

\begin{align}
\bar{X}_n=\frac{1}{n}\sum_{i=1}^{n}X_i.
\end{align}

Then, as $n\to\infty$, the standardized sample mean converges in
distribution to a standard normal random variable:

\begin{align}
Z_n=\frac{\bar{X}_n-\mu}{\sigma/\sqrt{n}}
\;\xrightarrow{d}\;
\mathcal{N}(0,1).
\end{align}

Equivalently, for large $n$,

\begin{align}
\bar{X}_n\ \approx\ \mathcal{N}\!\left(\mu,\ \frac{\sigma^{2}}{n}\right),
\qquad
\sum_{i=1}^{n}X_i\ \approx\ \mathcal{N}\!\left(n\mu,\ n\sigma^{2}\right).
\end{align}

\textbf{Relevance to Q12 (Binomial distribution):} a Binomial random
variable $X\sim B(N,p)$ is the sum of $N$ i.i.d.\ Bernoulli trials, each
with mean $p$ and variance $p(1-p)$. By the CLT, for large $N$,

\begin{align}
X\ \approx\ \mathcal{N}\!\left(Np,\ Np(1-p)\right),
\end{align}

which is why the Gaussian approximation to the Binomial distribution
uses exactly the mean $E(X)=Np$ and variance
$\operatorname{Var}(X)=Np(1-p)$ obtained in Section 12.

% =========================================================
\item Effect of Scalar Multiplication on a Determinant

For an $n\times n$ matrix $A$ and scalar $k$,

\begin{align}
\det(kA)=k^{n}\det(A). \label{det}
\end{align}
\end{enumerate}
%\end{document}
